%%%%%%%%%%%%%%%%%%%%%%%%%%%%%%%%%%%%%%%%%%%%%%%%%%%%%%%%%%%%%%%%%%%%%%%%%%%%%%%%
%2345678901234567890123456789012345678901234567890123456789012345678901234567890
%        1         2         3         4         5         6         7         8
 
\documentclass[letterpaper, 10 pt, conference]{ieeeconf}  % Comment this line out
                                                          % if you need a4paper
%\documentclass[a4paper, 10pt, conference]{ieeeconf}      % Use this line for a4
                                                           % paper

\IEEEoverridecommandlockouts                              % This command is only
                                                          % needed if you want to
                                                          % use the \thanks command
\overrideIEEEmargins
% See the \addtolength command later in the file to balance the column lengths
% on the last page of the document



% The following packages can be found on http:\\www.ctan.org
%\usepackage{graphics} % for pdf, bitmapped graphics files
%\usepackage{epsfig} % for postscript graphics files
%\usepackage{mathptmx} % assumes new font selection scheme installed
%\usepackage{times} % assumes new font selection scheme installed
%\usepackage{amsmath} % assumes amsmath package installed
%\usepackage{amssymb}  % assumes amsmath package installed

%\usepackage{overpic}
\usepackage{amsfonts}
%\usepackage{psfrag}
\usepackage{graphicx}
\usepackage[dvipsnames]{xcolor}
\usepackage{amsmath}
\usepackage{amssymb}
\usepackage{dsfont}

\usepackage{ifthen}
\usepackage{color}
\usepackage{cite}
%\input{steve.lib}

\usepackage{mathtools}
\usepackage{booktabs}
\usepackage{url}
\usepackage{hyperref}
% \usepackage{caption}
% \usepackage{subfig}
\usepackage[caption=false,font=footnotesize]{subfig}
%\usepackage{subfigure}
\usepackage{float}
\usepackage{multirow}
\usepackage{comment}
\usepackage{tikz}
\usepackage{pgfplots}
\pgfplotsset{compat=1.12}
%\usepackage{empheq}
\usepackage{algorithm}
\usepackage{algpseudocode}


\newcommand{\matt}[1]{\left[ \matrix{#1} \right]}
\def\scr#1{{\cal #1}}
\newcommand{\R}{{\rm I\!R}}
\def\eq#1{\begin{equation}#1\end{equation}}
\def\rep#1{(\ref{#1})}
\newcommand{\bbb}{\mathbb}
\newtheorem{theorem}{Theorem}
\newtheorem{lemma}{Lemma}
\newtheorem{remark}{Remark}
\newtheorem{assumption}{Assumption}
\newtheorem{definition}{Definition}
\newtheorem{proposition}{Proposition}
\newtheorem{corollary}{Corollary}
\newtheorem{conjecture}{Conjecture}
\newtheorem{observation}{Observation}
\newcommand{\dfb}{\stackrel{\Delta}{=}}
\def\qed{ \rule{.08in}{.08in}}
\DeclarePairedDelimiter\ceil{\lceil}{\rceil}
\DeclarePairedDelimiter\floor{\lfloor}{\rfloor}

\newcommand{\1}{\mathbf{1}}
\newcommand{\0}{\mathbf{0}}

\DeclareMathOperator*{\argmin}{arg\,min}
\DeclareMathOperator*{\argmax}{arg\,max}

%\newcommand{\matt}[1]{\left[ \matrix{#1} \right]}

%\newtheorem{definition}{Definition}
%\newtheorem{assumption}{Assumption}

\newcommand{\diag}{{\rm diag}}
\newcommand{\bestdelta}[1]{\textbf{\textcolor{BrickRed}{#1}}}
\newcommand{\bestwin}[1]{\textbf{\textcolor{NavyBlue}{#1}}}
\newcommand{\besttime}[1]{\textbf{\textcolor{ForestGreen}{#1}}}




\title{\LARGE \bf Graph Theory Guided Reinforcement Learning for Network Design with High Algebraic Connectivity under Vertex and Edge Constraints
%\thanks{Proofs of the assertions in this paper are available from the first author upon request and will appear in a full length version of this paper.
%}
}

%\author{ \parbox{3 in}{\centering Huibert Kwakernaak*
%         \thanks{*Use the $\backslash$thanks command to put information here}\\
%         Faculty of Electrical Engineering, Mathematics and Computer Science\\
%         University of Twente\\
%         7500 AE Enschede, The Netherlands\\
%         {\tt\small h.kwakernaak@autsubmit.com}}
%         \hspace*{ 0.5 in}
%         \parbox{3 in}{ \centering Pradeep Misra**
%         \thanks{**The footnote marks may be inserted manually}\\
%        Department of Electrical Engineering \\
%         Wright State University\\
%         Dayton, OH 45435, USA\\
%         {\tt\small pmisra@cs.wright.edu}}
%}

\author{Zeru Zhu \hspace{.3in} Marco Gamarra \hspace{.3in} Ji Liu
%\thanks{*The proofs of all assertions in this paper are omitted due to space limitations and can be found in the arXiv version of the paper \cite{arxivenergy}.}
\thanks{Z. Zhu is with the Department of Applied Mathematics and Statistics at Stony Brook University.
M. Gamarra is with the Air Force Research Laboratory.
J.~Liu is with the Department of Electrical and Computer Engineering at Stony Brook University.
%(ji.liu@stonybrook.edu).
Email addresses: 
\href{mailto:zeru.zhu@stonybrook.edu}{zeru.zhu@stonybrook.edu},
\href{mailto:marco.gamarra@us.af.mil}{marco.gamarra@us.af.mil},
\href{mailto:ji.liu@stonybrook.edu}{ji.liu@stonybrook.edu}.
}
%\thanks{The work of B.D.O.~Anderson is supported by Australian Research Council Discovery Projects DP160104500 and DP19010887 and by Data61-CSIRO.}
}



\begin{document}


\maketitle
\thispagestyle{empty}
\pagestyle{empty}


\begin{abstract}
This paper studies an open network design problem of constructing graph topologies with maximum algebraic connectivity under fixed numbers of vertices and edges. A scalable reinforcement learning approach is proposed, in which graph-theoretic insights guide the learning process and substantially reduce the search space. The method first selects and tailors a backbone structure based on graph density, using complete multipartite graphs in dense regimes and Cayley graphs in intermediate regimes, and then refines the remaining edges through a learning-based procedure. Numerical experiments demonstrate that the proposed approach is computationally efficient and consistently achieves high algebraic connectivity, outperforming existing methods in most cases.
\end{abstract}



\section{Introduction}


Over the past two decades, consensus has been extensively studied and has achieved notable success \cite{vicsekmodel,reza1,luc,ReBe05}, finding applications in a broad range of distributed control and computation problems \cite{fax,survey,nedic2009distributed,tacle,split}.
Consensus processes are fundamental to almost all cooperative multi-agent algorithms and, in many cases, significantly influence or even determine their convergence rate.


A typical continuous-time linear consensus process among $n$ agents over a connected simple graph $G$ can be modeled by a linear differential equation of the form $\dot x(t)=-Lx(t)$, where $x(t)$ is the state vector in $\R^n$ and $L$ is the ``Laplacian matrix'' of $G$. 
The $i$th entry of $x(t)$ represents the state of agent $i$. 
The multi-agent system is said to reach a consensus if all $x_i(t)$ are equal. 
For a simple graph $G$ with $n$ vertices, its Laplacian matrix is defined as $L=D-A$, where $D$ and $A$ denote its degree matrix and adjacency matrix, respectively. 
Here, $D$ is an $n\times n$ diagonal matrix whose $i$th diagonal entry equals the degree of vertex $i$, and $A$ is an $n\times n$ matrix whose $ij$th entry equals 1 if $(i,j)$ is an edge in $G$ and otherwise equals 0. 
It is well known that $L$ is positive semidefinite with smallest eigenvalue $0$, and its second smallest eigenvalue, called the algebraic connectivity of $G$, 
%and {\color{red}denoted by $a(\bbb G)$}, 
is positive if and only if $G$ is connected \cite{Fiedler73}. 
It has been shown that $\dot x(t)=-Lx(t)$ will reach a consensus if and only if $G$ is connected. 
More importantly, the (worst-case) convergence rate of consensus is characterized by the algebraic connectivity, with larger values leading to faster consensus \cite{reza1}. 

Motivated by the above discussion, a natural and fundamental problem is to design network topologies that achieve fast, or sometimes fastest, consensus. This problem has been extensively studied for a long time, among others \cite{moura,ming,bayen,dai,dong,ogiwara}. Despite these efforts, the following question remains largely open: Given fixed numbers of vertices and edges, which graphs maximize the algebraic connectivity?  
To date, this question has only been theoretically answered in certain special cases \cite{ogiwara,Shahbaz23}.
Beyond the challenges in theoretical analysis, the question is also computationally intractable in numerical experiments, especially for large-scale graphs. 
Indeed, it leads to an NP-hard combinatorial optimization problem \cite{MOSKAOYAMA2008677}.
Even though exhaustive combinatorial search is conceptually possible, its super-exponential growth makes it computationally impractical beyond very small graph sizes. The structural characterization of optimal graphs that maximize algebraic connectivity therefore remains poorly understood, let alone the associated computational complexity.


\subsection{Related Work and Benchmark Methods}

Due to the theoretical and computational hurdles, existing literature includes both heuristic methods, which are computationally practical but typically lack theoretical performance guarantees, and exact optimization formulations, such as MISDP-based approaches, whose main limitation is scalability.
It turns out that the problem of constructing graphs from scratch to achieve large algebraic connectivity, under constraints on the number of vertices and total edges, has received little attention.
The work in \cite{kurahashi2024depth} is restricted to the construction of special regular graphs. 
The work of \cite{Shahbaz23} proposes an algorithm centered around ``complete multipartite graphs'' to construct graphs for any feasible numbers of vertices and edges. 
A more general weighted version is studied in \cite{somisetty2025optimal}, where the problem is formulated as an equivalent MISDP, together with outer-approximation procedures for solving or bounding it.

Most existing works investigate a closely related augmentation problem: given a connected graph, how to add a fixed number of edges to maximally increase algebraic connectivity. This problem is NP-hard \cite{MOSKAOYAMA2008677}. Prior work includes SDP relaxations for Boolean edge-addition formulations \cite{ghosh2006growing}, as well as MISDP formulations \cite{somisetty2024spectral}. A number of heuristic and approximation methods have also been proposed. Notable among them are Fiedler vector based greedy methods \cite{ghosh2006growing,yin2024accelerating}, secular equation (for the Laplacian under a rank-one perturbation) guided edge selection \cite{kim2010bisection}, tabu search \cite{wei2014algebraic}, and Frank-Wolfe optimization with Fiedler vector gradients \cite{9981584}. In principle, these methods can be extended to construct a graph from scratch via different strategies, e.g., applying them iteratively.



%Lu and Liu, ``Fast Consensus Topology Design via Minimizing Laplacian Energy.''~\cite{lu2024fast} propose approximating the maximal algebraic connectivity by minimizing the ``Laplacian energy'' defined as follows.


It is worth emphasizing that all these existing methods are computationally expensive and therefore do not scale well to large-scale graph construction. 
%and providing no guarantees on their suboptimality relative to graphs with maximum algebraic connectivity.
We simulate representative methods in \cite{ghosh2006growing,8412478,wei2014algebraic,somisetty2024spectral,9981584,somisetty2025optimal} for constructing medium-sized graphs from scratch, and the results indicate that the Fiedler vector based greedy method \cite{ghosh2006growing} provides the best practical quality-runtime trade-off among full-construction baselines (see Subsection \ref{subsec:benchmark}).
For this reason, we adopt the Fiedler vector based greedy method as a benchmark for comparison with our proposed algorithm.


For a graph, any eigenvector corresponding to the second smallest eigenvalue of its Laplacian matrix is called a Fiedler vector of the graph \cite{Fiedler73}. We apply the Fiedler vector based greedy method to construct a graph from scratch by iteratively adding edges that maximize the increase in algebraic connectivity. Specifically, at each iteration, the increase caused by adding an edge $(i,j)$ is approximated via first-order eigenvalue perturbation as $(v_i - v_j)^2$, where $v$ is a selected Fiedler vector of the current graph; if $\lambda_2$ has multiplicity greater than one, $v$ is chosen from the corresponding eigenspace.

The method in \cite{Shahbaz23} constructs candidate graphs by specifying the complement as a union of complete components, which induces complete multipartite structures in the original graph and facilitates spectral optimization.

\subsection{Problem Formulation}
Let \(G=(V,E)\) be a connected simple undirected graph with \(|V|=n\) and \(|E|=m\).
%Its adjacency and degree matrices are denoted by \(A(G)\) and \(D(G)\), and its Laplacian matrix is \(L(G)=D(G)-A(G)\).
The Laplacian eigenvalues are ordered as \(0=\lambda_1(L)\le \lambda_2(L)\le \cdots\le \lambda_n(L)\).
We use \(\lambda_2(L(G))\) as the algebraic connectivity metric throughout the paper.
For fixed \(n\), define the connected graph family \(\mathcal G_n:=\{G:\ |V|=n,\ G\ \text{simple, undirected, connected}\}\).

For fixed \((n,m)\), define \(\mathcal G_{n,m}:=\{G\in\mathcal G_n:\ |E|=m\}\).
The design problem is
\begin{equation}
\max_{G\in\mathcal G_{n,m}}\ \lambda_2(L(G)).
\label{eq:main_opt}
\end{equation}
We refer to \eqref{eq:main_opt} as algebraic-connectivity maximization (ACM).
This is an NP-hard combinatorial optimization problem \cite{MOSKAOYAMA2008677}.

To compare graph budgets across different \(n\), we use the connected-range density
\begin{equation}
\rho
=
\frac{m-(n-1)}{\binom{n}{2}-(n-1)},
\qquad
\rho\in[0,1].
\label{eq:rho_connected}
\end{equation}
Thus, \(\rho=0\) corresponds to the minimum edge budget for connectivity (\(m=n-1\)), and \(\rho=1\) corresponds to the complete graph.
Given feasible \((n,\rho)\), we map back to \(m\) under the feasible range \(n-1\le m\le \binom n2\).



\subsection{Contributions}
Our main contributions are threefold.
First, we propose a density-dependent high-quality, low-cost backbone construction for fixed-\((n,m)\) ACM that combines complete multipartite and Cayley graph families, with a sample-wise parallelizable Cayley stage.
Second, we formulate RL completion as a residual edge-allocation stage initialized from the backbone, so learning focuses on the hard sequential part after structural initialization.
Third, across tested density levels and baselines, our backbone+RL method is the only non-exact approach that retains high \(\lambda_2\) quality while showing more favorable runtime scaling as graph size increases.


\section{Density-Dependent Backbone Construction}
Most existing methods in related work assume a given initial graph and then optimize edge additions.
In our setting, initialization itself is a design variable that strongly affects both search difficulty and final quality.
This motivates a density-dependent structural initialization that chooses the initial graph family according to edge density.

For each target \((n,m)\), we first construct a connected graph \(G_0=(V,E_0)\) with \(|E_0|=m_0\le m\), aiming for large \(\lambda_2(G_0)\) relative to its edge budget.
The remaining \(m-m_0\) edges are added in a later completion stage.
We refer to this initialized graph \(G_0\) as a \emph{backbone}.

\subsection{Envelope Backbone via Complete Multipartite Graphs}
For a fixed node count \(n\), each graph in \(\mathcal G_n\) corresponds to a discrete point in the \((\rho,\lambda_2)\) plane.
Since \(\lambda_2\) is monotone nondecreasing under edge addition, increasing edge density does not reduce the best achievable \(\lambda_2\);
hence, across the density spectrum, any maximizer of \eqref{eq:main_opt} for some feasible \(m\) must lie on the weak upper-left staircase, i.e., weak-maximal points under the order ``smaller \(\rho\), larger \(\lambda_2\).'' See Fig.~\ref{fig:staircase_random_n32}. We formalize this notion as follows.

\begin{figure}[t]
\centering
\includegraphics[width=0.95\linewidth]{plots/envelope/scatter_staircase_n32_s500_seed1234.png}
\caption{Random connected graphs (\(n=32\), 500 samples) in the \((\rho,\lambda_2)\) plane. Highlighted points illustrate the weak upper-left staircase (weak-maximal points).}
\label{fig:staircase_random_n32}
\end{figure}

\begin{definition}
Fix \(n\), and let \(\mathcal P\) be a finite set of points \(p=(\rho_p,\lambda_p)\) induced by a graph family in \(\mathcal G_n\).
For \(p,q\in\mathcal P\), we say \(p\) \emph{weakly dominates} \(q\) if \(\rho_p\le \rho_q\) and \(\lambda_p>\lambda_q\).
We say \(p\) \emph{strictly dominates} \(q\) if \(\rho_p\le \rho_q\), \(\lambda_p\ge \lambda_q\), and at least one inequality is strict.
A point is \emph{weak-maximal} (resp. \emph{strict-maximal}) if no other point in \(\mathcal P\) weakly dominates (resp. strictly dominates) it.
\end{definition}
This definition makes the geometric target explicit: weak-maximality characterizes where optimizers of \eqref{eq:main_opt} can lie, while strict-maximality enforces leftmost tie-breaking along horizontal levels.
Characterizing these maximal sets of \(\mathcal G_n\) directly is generally intractable.
We therefore study them within a structured subfamily where edge count and \(\lambda_2\) admit closed-form expressions.
We next introduce the complete multipartite family \(\mathcal K_n\), which is known to contain graphs with strong algebraic connectivity \cite{ogiwara,Shahbaz23}.
An illustrative subset of \(\mathcal K_n\) in the \((\rho,\lambda_2)\) plane is shown in Fig.~\ref{fig:envelope_staircase}.
\begin{definition}
Let \(r\in\{2,\dots,n\}\) and \(1\le n_1\le \cdots \le n_r\) with \(\sum_{i=1}^r n_i=n\).
The complete multipartite graph \(K_{n_1,\dots,n_r}\) has all edges between different parts and no edges within any part.
Denote by \(\mathcal K_n\subseteq\mathcal G_n\) the family of all such complete multipartite graphs on \(n\) vertices.
\end{definition}
It is known that when all parts have equal size (i.e., \(r\mid n\)), this graph is a maximizer of \(\lambda_2\) for the corresponding edge budget \cite{ishii2016extensions}.
We refer to this case as a balanced complete multipartite graph.

The first ingredient is a closed-form edge-count identity.
\begin{lemma}\label{lem:cmp_edge_count}
For \(G=K_{n_1,\dots,n_r}\),
\begin{equation}
m(G)=\sum_{1\le i<j\le r}n_i n_j
=\frac{1}{2}\left(n^2-\sum_{i=1}^r n_i^2\right).
\label{eq:cmp_edge}
\end{equation}
\end{lemma}
\begin{proof}
Edges appear only across different parts, so \(m(G)=\sum_{1\le i<j\le r} n_i n_j\).
Using \(\left(\sum_{i=1}^r n_i\right)^2=\sum_{i=1}^r n_i^2+2\sum_{1\le i<j\le r} n_i n_j\) and \(\sum_i n_i=n\), we obtain \eqref{eq:cmp_edge}.
\end{proof}
Hence edge minimization at fixed \(n\) is equivalent to maximizing \(\sum_i n_i^2\).

The second ingredient is a closed-form expression for \(\lambda_2\) in this family.
\begin{lemma}\label{lem:cmp_lambda2_formula}
Let \(G=K_{n_1,\dots,n_r}\) with \(2\le r\le n-1\) and \(n_1\le\cdots\le n_r\).
Then
\begin{equation}
\lambda_2(G)=n-n_r.
\label{eq:cmp_lambda2}
\end{equation}
For \(r=n\), \(G=K_n\) and \(\lambda_2(G)=n\).
\end{lemma}
\begin{proof}
Each vertex in part \(i\) has degree \(n-n_i\), so \(L(G)\) has diagonal blocks \((n-n_i)I_{n_i}\) and off-diagonal blocks \(-J_{n_i\times n_j}\), where \(J\) is the all-ones matrix.
For each \(i\), any vector supported on part \(i\) with zero coordinate sum is an eigenvector with eigenvalue \(n-n_i\), giving multiplicity \(n_i-1\).

Now consider the \(r\)-dimensional subspace of vectors that are constant on each part. It is invariant under \(L(G)\). The all-ones vector belongs to this subspace and gives eigenvalue \(0\). On the orthogonal complement of the all-ones vector in this subspace, a direct substitution shows \(Lx=nx\), so eigenvalue \(n\) has multiplicity \(r-1\).

Hence the spectrum consists of \(0\), \(n\) (mult. \(r-1\)), and \(n-n_i\) (mult. \(n_i-1\), \(i=1,\dots,r\)).
For \(2\le r\le n-1\), the smallest positive eigenvalue is \(n-n_r\), proving \eqref{eq:cmp_lambda2}.
For \(r=n\), all \(n_i=1\), so the \(n-n_i\) terms have zero multiplicity, and \(G=K_n\) gives \(\lambda_2(G)=n\).
\end{proof}
This shows that, within \(\mathcal K_n\), fixing \(\lambda_2=k\) is equivalent to fixing the largest part size as \(n-k\).

We now define the envelope point at a fixed \(\lambda_2\)-level.
\begin{definition}
Fix \(k\in\{1,\dots,n-2\}\).  
Among all \(G\in\mathcal K_n\) with \(\lambda_2(G)=k\), a graph with minimum edge count is called an envelope point at level \(k\).
\end{definition}
By definition, these are exactly strict-maximal points of \(\mathcal K_n\) in the \((\rho,\lambda_2)\) plane.

The next proposition gives an explicit construction and edge formula.
\begin{proposition}[Envelope Point at Level \(k\)]
\label{prop:env_formula}
Fix \(k\in\{1,\dots,n-2\}\), and write
\begin{equation}
n=q(n-k)+b,\qquad 0\le b<n-k.
\label{eq:env_division}
\end{equation}
An envelope point is obtained by taking \(q\) parts of size \(n-k\), and, if \(b>0\), one additional part of size \(b\).
Its edge count is
\begin{equation}
m_{\mathrm{env}}(k)=\frac{1}{2}\left(n^2-q(n-k)^2-b^2\right).
\label{eq:env_mk}
\end{equation}
\end{proposition}
\begin{proof}
By Lemma~\ref{lem:cmp_lambda2_formula}, fixing \(\lambda_2=k\) is equivalent to fixing the largest part size as \(n_r=n-k\).
Thus feasible partitions satisfy \(n_i\le n-k\) and \(\max_i n_i=n-k\).
By Lemma~\ref{lem:cmp_edge_count}, minimizing \(m(G)\) is equivalent to maximizing \(\sum_i n_i^2\).

In any maximizer of \(\sum_i n_i^2\) under these constraints, at most one part size is strictly smaller than \(n-k\).
Indeed, suppose there are two indices \(i\neq j\) with \(1\le n_i\le n_j<n-k\). Replacing \((n_i,n_j)\mapsto(n_i-1,n_j+1)\), where a zero part (if \(n_i=1\)) is discarded, keeps feasibility and strictly increases \(\sum_i n_i^2\), since \((n_i-1)^2+(n_j+1)^2-(n_i^2+n_j^2)=2(n_j-n_i)+2>0\).
This contradicts maximality.
Therefore all but at most one parts must have size exactly \(n-k\), giving the stated partition.
Writing \(n=q(n-k)+b\), \(0\le b<n-k\), yields that form, and substituting into \eqref{eq:cmp_edge} gives \eqref{eq:env_mk}.
\end{proof}

Proposition \ref{prop:env_formula} yields a closed-form staircase of strict-maximal points in \(\mathcal K_n\), providing a cheap dense-regime backbone library.
Detailed envelope-versus-baseline performance comparisons are deferred to Section~\ref{subsec:backbone_only}.

At very sparse densities, the first two envelope levels are key anchors: at \(k=1\), Proposition~\ref{prop:env_formula} gives \(K_{1,n-1}\), i.e., the star graph, which is optimal at the minimum connected budget \(m=n-1\) \cite{Mohar91}; at \(k=2\), it gives \(K_{2,n-2}\), whose optimality at the corresponding budget is a well-known open conjecture \cite{Kolokolnikov15}.

For backbone construction, given target \((n,m)\), choose \(k^\star=\max\{k\in\{1,\dots,n-2\}:m_{\mathrm{env}}(k)\le m\}\); because \(m_{\mathrm{env}}(k)\) is monotone nondecreasing, \(k^\star\) is found by integer binary search in \(O(\log n)\) time.
The envelope candidate is then the complete multipartite graph from Proposition \ref{prop:env_formula} at level \(k^\star\) (equivalently, using \eqref{eq:env_division}--\eqref{eq:env_mk}).
Its algebraic-connectivity score is \(\lambda_2=k^\star\) (or \(\lambda_2=n\) when \(m=\binom{n}{2}\)), and this candidate is passed to the routing/selection stage in Section III-C.
Algorithm~\ref{alg:env_backbone} summarizes this deterministic envelope construction.

\begin{algorithm}[t]
\caption{Envelope Backbone Construction}
\label{alg:env_backbone}
\begin{algorithmic}[1]
\State \textbf{Input:} \(n,m\)
\State \textbf{Output:} Envelope backbone \(G_{\mathrm{env}}\) with \(|E(G_{\mathrm{env}})|\le m\)
\If{\(m=\binom{n}{2}\)}
\State return \(K_n\) with \(\lambda_2=n\)
\EndIf
\State \(k^\star\leftarrow\max\{k\in\{1,\dots,n-2\}:m_{\mathrm{env}}(k)\le m\}\) by integer binary search on \eqref{eq:env_mk}
\State \(q\leftarrow \left\lfloor n/(n-k^\star)\right\rfloor\), \(b\leftarrow n-q(n-k^\star)\)
\State build \(K_{n_1,\dots,n_r}\) with \(q\) parts of size \(n-k^\star\) and one part \(b\) if \(b>0\), \(G_{\mathrm{env}}\leftarrow K_{n_1,\dots,n_r}\)
\State return \(G_{\mathrm{env}}\) and \(\lambda_2(G_{\mathrm{env}})=k^\star\)
\end{algorithmic}
\end{algorithm}

\begin{figure}[t]
    \centering
    \includegraphics[width=0.95\linewidth]{plots/envelope/multipartite_scatter_staircase_n36_k2_3_4_5.png}
    \caption{Scatter of complete 2-, 3-, 4-, and 5-partite graphs in \(\mathcal K_{36}\) over the \((\rho,\lambda_2)\) plane.}
    \label{fig:envelope_staircase}
\end{figure}

\subsection{Cayley Backbone via Coset-Based Generating Sets}
While envelope points from \(\mathcal K_n\) are strong in dense regimes, our experiments show a clear weakness for \(\rho<\tfrac12\), where their \(\lambda_2\) is non-competitive (see Section~\ref{subsec:backbone_only}).
Therefore, the need for a second backbone family is driven not only by limited envelope coverage across feasible density levels, but also by quality degradation in low-to-intermediate densities.
Balanced complete multipartite graphs are a natural anchor because they are guaranteed optimal at their corresponding budgets \cite{ishii2016extensions}.
This raises a natural question: beyond \(\mathcal K_n\), is there a structured graph family that extends balanced complete multipartite graphs, yet can realize stronger candidates in other density ranges?
Cayley graphs provide such a generalization; as shown below, complete balanced multipartite graphs arise as a special Cayley case.

We next give a formal definition of Cayley graphs.
\begin{definition}
Let \(\Gamma\) be a finite group and \(S\subseteq \Gamma\setminus\{e\}\) with \(S=S^{-1}\).
The undirected Cayley graph \(\mathrm{Cay}(\Gamma,S)\) has vertex set \(\Gamma\) and edge set \(\{\{g,gs\}: g\in\Gamma,\ s\in S\}\).
\end{definition}
The inverse-closed condition ensures undirected edges, while group actions provide strong symmetry and reusable structure for sampling.

The next proposition establishes that the balanced complete multipartite case is a special instance of Cayley graphs.
\begin{proposition}[Balanced Multipartite as Cayley]
\label{prop:balanced_cmp_cayley}
Let \(\Gamma\) be a finite group, and let \(H<\Gamma\) be a proper subgroup of index \(r\ge 2\).
Define \(S=\Gamma\setminus H\).
Then \(\mathrm{Cay}(\Gamma,S)\) is a complete balanced \(r\)-partite graph whose parts are the left cosets of \(H\).
\end{proposition}
\begin{proof}
Since \(e\in H\), we have \(e\notin S\). Also \(S^{-1}=S\), so \(\mathrm{Cay}(\Gamma,S)\) is an undirected Cayley graph.
Let \(\Gamma=\bigsqcup_{\ell=1}^r g_\ell H\) be the left-coset decomposition. Each part has size \(|H|\), so the partition is balanced.
For \(x,y\in\Gamma\), if they are in the same coset then \(x^{-1}y\in H\), hence \(x^{-1}y\notin S\), so no edge.
If they are in different cosets then \(x^{-1}y\notin H\), hence \(x^{-1}y\in S\), so an edge exists.
Therefore adjacency is exactly across distinct cosets, i.e., a complete balanced \(r\)-partite graph.
\end{proof}
Proposition~\ref{prop:balanced_cmp_cayley} gives the key motivation for our construction: index-\(r\) coset-based generators recover the balanced \(r\)-partite endpoint when all nontrivial cosets are selected, and random inverse-closed sampling from these cosets extends that structure to lower density levels.
Numerically, we observe that generic random Cayley candidates are already competitive, while coset-based Cayley candidates are typically stronger.

An index-\(r\) subgroup requires \(r\mid n\), so direct coset construction is not always available.
To handle arbitrary \(n\), we use a lift-and-delete extension: construct a Cayley graph on \(n+j\) vertices where divisibility holds, then delete \(j\) vertices to return to size \(n\).

\begin{lemma}[Vertex-Transitivity of Cayley Graphs]
\label{lem:cayley_vertex_transitive}
Let \(\Gamma\) be a finite group and \(S\subseteq \Gamma\setminus\{e\}\) with \(S=S^{-1}\).
Then \(\mathrm{Cay}(\Gamma,S)\) is vertex-transitive.
\end{lemma}
\begin{proof}
This is a standard property of Cayley graphs; see Theorem 3.1.2 in \cite{godsil_royle_agt_2001}.
\end{proof}
By Lemma~\ref{lem:cayley_vertex_transitive}, deleting one vertex is equivalent to deleting any other vertex up to graph isomorphism, which supports the one-step lift-and-delete operation.
Empirically, this extension remains competitive after projection back to \(n\).

How the choice of group family and subgroup affects algebraic connectivity has not been systematically investigated and is left as a future direction.
For implementation simplicity, we use standard group families (e.g., dihedral or cyclic) and fixed subgroup templates.
We therefore define the Cayley backbone as a coset-based generating-set Cayley construction, with lift-and-delete applied when divisibility constraints prevent direct index-\(r\) construction.
Algorithm~\ref{alg:cayley_backbone} gives the corresponding fixed-index-\(r\) backbone procedure.
When the feasible degree is too small (e.g., \(k=1\) or \(k=2\)), sampled Cayley candidates can be disconnected.

\begin{algorithm}[t]
\caption{Index-\(r\) Cayley Backbone Construction}
\label{alg:cayley_backbone}
\begin{algorithmic}[1]
\State \textbf{Input:} \(n,m,r\), sample count \(t\ge 1\)
\State \textbf{Output:} Cayley backbone \(G_{\mathrm{cay}}\)
\State compute lift-delete \(j\) and \(k_{\max}\leftarrow\left\lfloor\frac{2m}{n-j}\right\rfloor\) for fixed \(r\)
\State choose preset group \(\Gamma\) and index-\(r\) subgroup \(H\) (with lift if needed)
\State set \(n_{\mathrm{lift}}\leftarrow n+j\), \(k\leftarrow k_{\max}\)
\If{\(k\) is not inverse-closed feasible}
\State set \(k\leftarrow \max\{\tilde{k}\le k:\ \tilde{k}\text{ is inverse-closed feasible}\}\)
\EndIf
\State set \(\widehat{\tau}\leftarrow \frac{n_{\mathrm{lift}}^2}{2}\cdot\frac{r-1}{r}\), \(\widehat{m}\leftarrow m+jk\)
\If{\(j=0\) and \(m=\widehat{\tau}\)}
\State return balanced complete \(r\)-partite graph
\EndIf
\State initialize candidate set \(\mathcal{C}\leftarrow\emptyset\)
\If{\(\widehat{m}\le \widehat{\tau}\)}
\State sample \(t\) inverse-closed generating sets \(S_1,\ldots,S_t\) of size \(k\) from \(\Gamma\setminus H\)
\Else
\State for each \(s=1,\ldots,t\), form \(S_s\) by including all elements in \(\Gamma\setminus H\), then sampling the remaining \(k-\lvert\Gamma\setminus H\rvert\) generators (inverse-closed) from \(H\setminus\{e\}\)
\EndIf
\State for each sampled \(S_s\), build \(G_s\leftarrow \mathrm{Cay}(\Gamma,S_s)\) on \(n_{\mathrm{lift}}\) vertices
\State delete \(j\) vertices from \(G_s\) to obtain \(n\)-vertex candidate \(\widetilde{G}_s\), then update \(\mathcal{C}\leftarrow \mathcal{C}\cup\{\widetilde{G}_s\}\)
\State return \(G_{\mathrm{cay}}=\arg\max_{G\in\mathcal{C}}\lambda_2(G)\)
\end{algorithmic}
\end{algorithm}

At the coverage level, index-\(r\) Cayley construction reaches edge ratios up to approximately \((r-1)/r\). Hence indices \(2,3,4\) cover low-to-intermediate densities up to approximately \(3/4\), i.e., most regimes where envelope points are not dominant in practice.
Based on this coverage property, we now instantiate a deterministic density-routing rule that decides which family (or both) is evaluated at each target density, and how Cayley indices are selected under edge feasibility.

\subsection{Density Routing and Backbone Selection Rule}
We operationalize the above family-level observations using a compact overlap rule:
envelope is used for \(\rho\in[0,0.1)\cup[0.5,1]\), Cayley is used for \(\rho\in[0.05,0.75]\), and both families are evaluated on the overlap windows \([0.05,0.1)\) and \([0.5,0.75]\).
For Cayley candidates, we use index-\(2/3/4\) coset constructions with inverse-closed sampling under edge-budget feasibility.
Specifically, index-2 uses dihedral/semidirect templates, while indices 3 and 4 use cyclic templates.
For each index \(r\in\{2,3,4\}\), let \(j_r\) be the lift-delete count.
We define \(k_{\max}^{(r)}\) as the largest budget-feasible Cayley degree after lift-delete adjustment, \(m_{\mathrm{post}}^{(r)}\) as the corresponding post-delete edge count, and \(d_{\mathrm{post}}^{(r)}\) as the density-based pocket coordinate:
for simplicity in this pocket computation, we assume the deleted vertices have no edges among themselves.
\begin{equation}
\begin{aligned}
k_{\max}^{(r)}&=\left\lfloor\frac{2m}{n-j_r}\right\rfloor,\quad
m_{\mathrm{post}}^{(r)}=\left\lfloor\frac{k_{\max}^{(r)}(n-j_r)}{2}\right\rfloor,\\
d_{\mathrm{post}}^{(r)}&=\frac{2m_{\mathrm{post}}^{(r)}}{n^2}.
\end{aligned}
\label{eq:cayley_pocket}
\end{equation}
\(d_{\mathrm{post}}^{(r)}\) is the post-delete average-degree-to-\(n\) ratio; at balanced index-\(r\) endpoints it equals \(1-\frac{1}{r}\), i.e., \(1/2,2/3,3/4\) for \(r=2,3,4\).
Density-based pocket tests are then applied as \(d_{\mathrm{post}}^{(2)}\in(0,\tfrac{1}{2})\), \(d_{\mathrm{post}}^{(3)}\in(\tfrac{1}{2},\tfrac{2}{3})\), and \(d_{\mathrm{post}}^{(4)}\in(\tfrac{2}{3},\tfrac{3}{4})\), with the smallest valid index selected when multiple are valid.
This matches the practical coverage of index-\(2/3/4\) Cayley constructions at approximately \(1/2,2/3,3/4\), respectively (exact in no-lift cases). Empirically, for \(\rho>3/4\), envelope points are typically stronger (Section~\ref{subsec:backbone_only}), so routing favors envelope there.
In overlap windows, both families are evaluated; each attempts to produce a backbone, and when both are feasible, the final family choice is made after completion by comparing final \(\lambda_2\).
In Cayley-only windows, if sampled Cayley candidates are infeasible/disconnected (typically at very small feasible degree), envelope is used as fallback.

\subsection{Backbone Construction Algorithm and Scalability}
The routing rule in Section III-C defines \emph{when} each family is used; Algorithm~\ref{alg:backbone} consolidates \emph{how} the backbone stage is executed end-to-end by combining envelope construction (Algorithm~\ref{alg:env_backbone}), Cayley construction (Algorithm~\ref{alg:cayley_backbone}), overlap handling, and fallback.

For fixed \((n,m)\), envelope-level lookup is \(O(\log n)\) via monotone integer binary search. In the Cayley branch, \(t\) sampled candidates are constructed and evaluated; these candidate construction/evaluation tasks are independent across samples and can be parallelized. By contrast, FVG-style completion is path-dependent and sequential.

In practice, we use fixed \(t=500\) by default. Increasing \(t\) beyond 500 (e.g., \(t=5000\)) gives only marginal quality gains in our tested ranges, and runtime is quickly dominated by per-candidate \(\lambda_2\) evaluation rather than candidate construction. Although one would intuitively expect required \(t\) to increase with \(n\), fixed-\(t\) performance does not degrade rapidly as \(n\) grows in our experiments.

\begin{algorithm}[t]
\caption{Density-Dependent Backbone Construction}
\label{alg:backbone}
\begin{algorithmic}[1]
\State \textbf{Input:} \(n,m\), Cayley sampling count \(t\)
\State \textbf{Output:} Backbone candidate list \(\mathcal B\)
\State compute \(\rho\) from \((n,m)\) via \eqref{eq:rho_connected}
\State set \(\texttt{env\_active}\leftarrow[\rho\in[0,0.1)\cup[0.5,1]]\), \(\texttt{cay\_active}\leftarrow[\rho\in[0.05,0.75]]\)
\State initialize \(\mathcal B\leftarrow[\,]\)
\If{\(\texttt{env\_active}\)}
\State \(G_{\mathrm{env}}\leftarrow\) Algorithm~\ref{alg:env_backbone}\((n,m)\)
\State append \(G_{\mathrm{env}}\) to \(\mathcal B\)
\EndIf
\If{\(\texttt{cay\_active}\)}
\State select index \(r\) by \eqref{eq:cayley_pocket}, and run Algorithm~\ref{alg:cayley_backbone}\((n,m,r,t)\)
\If{a connected Cayley candidate exists}
\State append \(G_{\mathrm{cay}}\) to \(\mathcal B\)
\EndIf
\EndIf
\If{\(\mathcal B\) is empty}
\State \(G_{\mathrm{env}}\leftarrow\) Algorithm~\ref{alg:env_backbone}\((n,m)\) \Comment{fallback}
\State append \(G_{\mathrm{env}}\) to \(\mathcal B\)
\EndIf
\State return \(\mathcal B\)
\end{algorithmic}
\end{algorithm}

\section{RL Completion}
After backbone construction, we solve only the residual edge-allocation problem by a sequential policy.
Given an initial connected backbone \(G_0=(V,E_0)\) with \(|E_0|\le m_{\mathrm{target}}\), each step adds one non-edge for \(T=m_{\mathrm{target}}-|E_0|\) steps.
Thus the RL stage does not design a graph from scratch; it refines a density-guided backbone under a fixed residual budget.

\subsection{Residual Edge-Addition MDP}
A state is the current graph \(G_t=(V,E_t)\), and an action is a non-edge \((i,j)\notin E_t\).
The transition is \(E_{t+1}=E_t\cup\{(i,j)\}\), and the episode terminates at \(t=T\) (equivalently \(|E_t|=m_{\mathrm{target}}\)).
In our implementation, episodes are sampled over a curriculum of graph sizes and random target densities, then mapped to connected-range edge budgets by \eqref{eq:rho_connected}.

\subsection{State and Candidate Representation}
At each step, we first form a candidate action pool by keeping only the top-\(k\) non-edges ranked by the Fiedler score \((\phi_2(i)-\phi_2(j))^2\) (\(k=128\)), where \(\phi_2,\phi_3\) are eigenvectors of \(\lambda_2,\lambda_3\).
For each candidate, the policy uses node, pair, and global features:
\[
\text{node: }\big[d_i/(n-1),\,d_i^{(2)}/(n-1),\,\phi_2(i),\,\phi_3(i),\,C_i\big],
\]
\[
\text{pair: }\left[
\begin{aligned}
&(\phi_2(i)-\phi_2(j))^2,\;|\phi_3(i)-\phi_3(j)|,\\
&\mathrm{CN}^{\mathrm{norm}},\;\mathrm{Jaccard},\;
\min\{\mathrm{dist}(i,j),d_{\mathrm{cap}}\}/d_{\mathrm{cap}}
\end{aligned}
\right],
\]
\[
\text{global: }\big[n,\,\rho_{\mathrm{target}},\,\rho_t,\,\lambda_2/n,\,\lambda_3/n\big].
\]
Here \(d_i\) is node degree, \(d_i^{(2)}\) is the two-hop neighbor count (excluding one-hop neighbors), \(C_i\) is local clustering coefficient, \(\mathrm{CN}^{\mathrm{norm}}=|N(i)\cap N(j)|/(n-2)\), \(\mathrm{Jaccard}=|N(i)\cap N(j)|/|N(i)\cup N(j)|\), and \(\mathrm{dist}(i,j)\) is shortest-path distance in \(G_t\).

\subsection{Policy and Value Networks}
Let \(x_i\) denote node features and \(h_i\) denote the corresponding GAT embedding.
For each candidate non-edge \((i,j)\), a symmetric edge representation is built from
\[
|h_i-h_j|,\quad h_i\odot h_j,\quad |x_i-x_j|,\quad x_i\odot x_j,
\]
then concatenated with pair and global features and scored by an MLP.
The action distribution is the softmax over candidate scores.
A value head shares the GAT trunk and predicts \(V(s_t)\) from pooled node embedding plus global features.

\subsection{Training Objective}
Completion is trained with PPO, using a value baseline and entropy regularization.
The per-step reward is
\[
r_t=\frac{\lambda_2(G_{t+1})-\lambda_2(G_t)}{n}
\;+\;
\beta\cdot \mathbf{1}_{\mathrm{terminal}}\frac{\lambda_2(G_T)}{n},
\]
where \(\beta\) is the terminal-bonus coefficient.
Returns are Monte Carlo discounted sums,
\[
A_t=\mathcal{R}_t-V_\psi(s_t), \qquad
\mathcal{R}_t=\sum_{\tau=t}^{T-1}\gamma^{\tau-t}r_\tau .
\]
With importance ratio \(r_t(\theta)=\pi_\theta(a_t|s_t)/\pi_{\theta_{\mathrm{old}}}(a_t|s_t)\), the clipped policy term is
\[
\mathcal L_{\mathrm{PPO}}^{\pi}
=
-\mathbb E\!\left[
\min\!\left(
r_t(\theta)A_t,\;
\mathrm{clip}\!\left(r_t(\theta),1-\epsilon,1+\epsilon\right)A_t
\right)\right],
\]
and the full objective is
\[
\mathcal L_{\mathrm{PPO}}
=
\mathcal L_{\mathrm{PPO}}^{\pi}
+c_v\,\mathbb E\!\left[(V_\psi(s_t)-\mathcal{R}_t)^2\right]
-c_e\,\mathbb E\!\left[\mathcal H(\pi_\theta(\cdot|s_t))\right].
\]
PPO is trained for \(1.2\times10^6\) environment steps using a three-phase size curriculum, with \(4\times10^5\) steps per phase over \(n\in[12,16]\), \(n\in[16,32]\), and \(n\in[16,64]\).

\section{Numerical Results}
\subsection{Experimental Setup}
All methods are evaluated under the same edge-budget mapping \eqref{eq:rho_connected}.
For each tested \(n\), we sweep 100 density points over \(\rho\in[0,1]\), and each point maps to the same \((n,m)\) target across methods.
The quality metric is final \(\lambda_2\) at the target budget, and runtime is wall-clock time on identical hardware.
Detailed baseline hyperparameter settings are listed in Appendix~\ref{app:baseline_hparams}.

\subsection{Baseline Comparison Among Existing Methods}\label{subsec:benchmark}
This subsection benchmarks full-construction baselines.
Compared methods include constructive heuristics (FVG \cite{ghosh2006growing}, MDMD \cite{8412478}), local search (WTS \cite{wei2014algebraic}, k-opt \cite{somisetty2024spectral}), relaxation/optimization (MAC \cite{9981584}, SDP-step \cite{wei2014algebraic}), and exact OA-PMC \cite{somisetty2025optimal}. Tables \ref{tab:baseline_lambda2_summary}--\ref{tab:baseline_runtime_summary} summarize quality/runtime, and Figures \ref{fig:baseline_lambda2}--\ref{fig:baseline_lambda2_small} show representative quality curves.

\begin{figure}[t]
    \centering
    \includegraphics[width=0.95\linewidth]{plots/baseline/lambda2_n16_rho_mean.png}
    \caption{Baseline quality comparison at \(n=16\) (\(\lambda_2\) over \(\rho\)).}
    \label{fig:baseline_lambda2}
\end{figure}

\begin{figure}[t]
    \centering
    \includegraphics[width=0.95\linewidth]{plots/baseline/lambda2_n64_rho_mean.png}
    \caption{Baseline quality comparison at \(n=64\) (\(\lambda_2\) over \(\rho\)).}
    \label{fig:baseline_lambda2_small}
\end{figure}

\begin{table}[t]
    \centering
    \caption{Quality summary vs FVG (\(\Delta\lambda_2\%\) / win\% / match\%). Red bold: best \(\Delta\lambda_2\); blue bold: best win\%.}
    \label{tab:baseline_lambda2_summary}
    \footnotesize
    \setlength{\tabcolsep}{1.3pt}
    \renewcommand{\arraystretch}{1.05}
    \resizebox{\columnwidth}{!}{%
    \begin{tabular}{lccccc}
        \toprule
        Method & \(n=8\) & \(n=16\) & \(n=32\) & \(n=64\) & \(n=128\) \\
        \midrule
        MAC      & -19.0/0.0/20.8  & -25.0/2.1/2.1  & -33.4/0.0/0.0  & -37.3/0.0/0.0  & -40.9/0.0/0.0 \\
        MDMD     & -1.0/10.4/64.6  & -9.2/4.2/10.4  & -9.5/2.1/2.1   & -9.0/0.0/0.0   & -8.5/0.0/0.0 \\
        FVG      & 0.0/0.0/100.0   & 0.0/0.0/100.0  & 0.0/0.0/100.0  & 0.0/0.0/100.0  & 0.0/0.0/100.0 \\
        WTS      & 11.4/56.3/43.8  & 6.4/\bestwin{88.2}/7.6   & -0.1/16.0/0.0  & -3.0/1.4/0.0   & -3.5/0.0/0.0 \\
        k-opt    & 8.7/52.1/47.9   & 4.9/77.1/22.9  & 1.6/\bestwin{89.6}/10.4  & 0.4/81.3/18.8  & 0.1/65.3/1.4 \\
        SDP-step & 3.3/25.0/52.1   & 1.7/52.1/8.3   & 1.8/70.8/2.1   & --             & -- \\
        OA-PMC   & \bestdelta{16.6}/\bestwin{65.0}/35.0  & --             & --             & --             & -- \\
        BB-RL(Ours) & 14.1/48.9/40.4 & \bestdelta{20.8}/67.3/9.2 & \bestdelta{14.5}/81.6/4.1  & \bestdelta{7.4}/\bestwin{81.6}/2.0   & \bestdelta{3.3}/\bestwin{71.4}/1.0 \\
        \bottomrule
    \end{tabular}
    }
\end{table}

\begin{table}[t]
    \centering
    \caption{Runtime summary (mean seconds). Green bold: fastest method at each \(n\).}
    \label{tab:baseline_runtime_summary}
    \footnotesize
    \setlength{\tabcolsep}{1.3pt}
    \renewcommand{\arraystretch}{1.05}
    \resizebox{\columnwidth}{!}{%
    \begin{tabular}{lccccc}
        \toprule
        Method & \(n=8\) & \(n=16\) & \(n=32\) & \(n=64\) & \(n=128\) \\
        \midrule
        MAC      & \(3.06\times10^{-3}\) & \besttime{\(4.70\times10^{-3}\)} & \besttime{\(3.47\times10^{-2}\)} & \besttime{\(9.37\times10^{-2}\)} & \besttime{\(3.54\times10^{-1}\)} \\
        MDMD     & \besttime{\(6.64\times10^{-4}\)} & \(6.06\times10^{-3}\) & \(4.14\times10^{-2}\) & \(4.04\times10^{-1}\) & \(4.75\) \\
        FVG      & \(2.05\times10^{-3}\) & \(7.80\times10^{-3}\) & \(1.87\times10^{-1}\) & \(1.77\) & \(24.73\) \\
        WTS      & \(5.60\times10^{-1}\) & \(9.32\times10^{-1}\) & \(2.05\) & \(5.93\) & \(33.72\) \\
        k-opt    & \(5.20\times10^{-2}\) & \(3.60\times10^{-1}\) & \(1.94\) & \(6.91\) & \(43.65\) \\
        SDP-step & \(5.29\times10^{-1}\) & \(4.54\) & \(85.87\) & -- & -- \\
        OA-PMC   & \(64.97\) & -- & -- & -- & -- \\
        BB-RL(Ours) & \(7.68\times10^{-3}\) & \(2.46\times10^{-2}\) & \(7.55\times10^{-2}\) & \(2.83\times10^{-1}\) & \(1.67\) \\
        \bottomrule
    \end{tabular}
    }
\end{table}

Taken together, these results show a clear quality-runtime hierarchy: MAC/MDMD are fast but weaker than FVG, WTS/k-opt are usually slower and only occasionally better at small \(n\), and SDP-step/OA-PMC are stronger (or equal) when applicable but much more expensive. Therefore, FVG is used as the primary practical full-construction baseline.

\subsection{Backbone-Only Performance Across Density Regimes}\label{subsec:backbone_only}
To isolate structural contribution, we compare backbone outputs across methods (without RL completion).
We include EIG \cite{Shahbaz23} as an additional complete-multipartite structural baseline.
Dense-regime quality is shown in Figures \ref{fig:backbone_only_lambda2}--\ref{fig:backbone_only_lambda2_dense}, while low-to-intermediate density quality is shown in Figure \ref{fig:backbone_only_lambda2_low}.
The optimality-guaranteed subset among envelope points is justified by the criterion in Corollary 4.6 of \cite{Shahbaz23}.

\begin{figure}[t]
    \centering
    \includegraphics[width=0.95\linewidth]{plots/envelope/compare_n60.png}
    \caption{Envelope vs FVG and EIG at \(n=60\).}
    \label{fig:backbone_only_lambda2}
\end{figure}

\begin{figure}[t]
    \centering
    \includegraphics[width=0.95\linewidth]{plots/envelope/compare_n512_zoomin.png}
    \caption{Envelope vs FVG and EIG at \(n=512\) and \(\rho>0.9\).}
    \label{fig:backbone_only_lambda2_dense}
\end{figure}

\begin{figure}[t]
    \centering
    \includegraphics[width=0.95\linewidth]{plots/cayley/compare_n64_density_0_0p6.png}
    \caption{Low-intermediate quality comparison at \(n=64\).}
    \label{fig:backbone_only_lambda2_low}
\end{figure}

Taken together, these results show that envelope dominates the dense regime, Cayley provides competitive low-to-intermediate candidates, and backbone construction captures regime-dependent structure before completion.

\subsection{Backbone+RL Performance}
We report head-to-head comparison on identical \((n,m)\) targets: backbone-initialized RL and FVG. Figures \ref{fig:backbone_rl_lambda2_n47} and \ref{fig:backbone_rl_lambda2_n64}, together with Tables \ref{tab:baseline_lambda2_summary} and \ref{tab:baseline_runtime_summary}, summarize \(\lambda_2\) and runtime.
Across densities, backbone+RL is the only non-exact tested method that avoids the usual quality--scalability trade-off: as \(n\) grows, it maintains high \(\lambda_2\) while runtime scales more favorably than FVG. Reported runtimes are serial; independent Cayley samples admit straightforward parallel speedup, unlike sequential FVG.

\begin{figure}[t]
    \centering
    \includegraphics[width=0.95\linewidth]{plots/RL/lambda2_n47_backbone_rl_vs_fvg_markers.png}
    \caption{\(\lambda_2\) vs density at \(n=47\): backbone+RL and FVG.}
    \label{fig:backbone_rl_lambda2_n47}
\end{figure}

\begin{figure}[t]
    \centering
    \includegraphics[width=0.95\linewidth]{plots/RL/lambda2_n64_backbone_rl_vs_fvg_markers.png}
    \caption{\(\lambda_2\) vs density at \(n=64\): backbone+RL and FVG.}
    \label{fig:backbone_rl_lambda2_n64}
\end{figure}

\section{Conclusion}
We presented a density-dependent framework for fixed-\((n,m)\) ACM that separates structural initialization from completion.
A high-quality low-cost backbone plus RL completion gives the best non-exact quality/runtime scaling among tested baselines as \(n\) increases.
As \(n\) grows, cheaper heuristics lose quality while stronger optimization-based methods become increasingly expensive, supporting backbone-first completion as a practical direction for scalable ACM.

\appendix
\section{Baseline Hyperparameters}
\label{app:baseline_hparams}
Unless otherwise stated, baseline methods use fixed settings across all reported runs. WTS (tabu search) uses random initialization, \(1000\) iterations, tabu tenure \(100\), and neighborhood controls with at most \(8\) old edges per iteration, \(3\) sampled neighbors per old edge, and \(1\) random jump per old edge. MAC uses Frank--Wolfe iterations \(50\), duality-gap tolerance \(10^{-6}\), and Fiedler-topk initialization. k-opt uses \((k,m)=(2,10)\), maximum rounds \(10\), and add/delete combination caps \(500/500\). SDP-step uses MOSEK with maximum iterations \(5\times 10^4\) and tolerance \(10^{-6}\). OA-PMC uses GUROBI with maximum iterations \(60\) and tolerance \(10^{-6}\).

\bibliographystyle{unsrt}
\bibliography{susie,jicareer,consensus,push,related_work}


\end{document}
